% Source: cft.pdf, pages 1-10 (typed handout). Title on the source: "Phys. 541, Notes on Conformal
% Invariance and Continuous Phase Transitions in 2D". The notes carry no lecture number.
\subsection{Conformal transformations and the Virasoro algebra}

These are some brief notes on the use of conformal invariance in the study of continuous phase transitions in two
dimensions. At the end, there is a list of a few references that give much more detailed discussions.

A continuous (second-order) phase transition is generically accompanied by an infinite correlation length, and hence
scale invariance, at the critical point. In $d = 2$ dimensions, this invariance is enhanced to conformal invariance.
These transformations preserve angles between vectors in the 2D space. It is convenient to represent $(x,y)$ in 2D as
a complex variable, $z = x+\ii y$. A conformal mapping $f(z)$ is generically a function that is analytic. A conformal
transformation may thus be written as a Taylor--Laurent series expansion
\begin{equation}\label{eq:cft-laurent}
  f(z) = \sum_{n=-\infty}^{\infty} a_nz^n .
\end{equation}
In a Taylor expansion, the coefficients are obtained as derivatives, which leads to the consideration of the
differential operators
\begin{equation}\label{eq:cft-Ln}
  L_n = -z^{n+1}\frac{\dd}{\dd z} .
\end{equation}
(The minus sign is included by convention to yield the algebra \eqref{eq:cft-witt} below.) The operator
$L_0 = -z\,\dd/\dd z = -\dd/\dd\ln z$. With $z = re^{\ii\theta}$, it follows that if one holds $\theta$ constant, then
$L_0$ generates dilatations. Consequently, the eigenvalues of $L_0$ are the (conformal) scaling dimensions, denoted
generically as $h$. The set of $L_n$ generate an algebra, as follows:
\begin{align}
  L_mL_n &= z^{m+1}\frac{\dd}{\dd z}\left(z^{n+1}\frac{\dd}{\dd z}\right) \nonumber\\
         &= (n+1)z^{m+n+1}\frac{\dd}{\dd z} + z^{m+n+2}\frac{\dd^2}{\dd z^2}
\end{align}
and
\begin{align}
  L_nL_m &= z^{n+1}\frac{\dd}{\dd z}\left(z^{m+1}\frac{\dd}{\dd z}\right) \nonumber\\
         &= (m+1)z^{n+m+1}\frac{\dd}{\dd z} + z^{n+m+2}\frac{\dd^2}{\dd z^2} .
\end{align}
Hence, for the commutator,
\begin{align}
  [L_m,L_n] \equiv L_mL_n - L_nL_m &= [(n+1)-(m+1)]z^{m+n+1}\frac{\dd}{\dd z} \nonumber\\
  &= (n-m)z^{m+n+1}\frac{\dd}{\dd z} = (m-n)L_{m+n} . \label{eq:cft-witt}
\end{align}

\subsubsection*{The Witt and Virasoro algebras}

Eq.~\eqref{eq:cft-witt} is the classical form of the conformal algebra and is called the Witt algebra. One can
construct a central extension term, which defines the full conformal algebra and is commonly called the Virasoro
algebra:
\begin{equation}\label{eq:cft-virasoro}
  [L_m,L_n] = (m-n)L_{m+n} + \frac{c}{12}m(m^2-1)\delta_{m,-n} .
\end{equation}
Here, $c$ is called the central charge, and the factor of $1/12$ is a convention. For a discussion, see, e.g., Green,
Schwarz, and Witten, \emph{Superstring Theory}, vol.~1 pp.~79--81. One can also define a corresponding Virasoro algebra
involving generators $\bar L_n = -\bar z^{n+1}\,\dd/\dd\bar z$, where $\bar z = z^* = x-\ii y$.

The analysis of the Virasoro algebra has some aspects in common with the analysis of Lie algebras, but there are
important differences. One important difference is immediately apparent; since all positive and negative values of
$n$ enter in the Taylor--Laurent series, it follows that, in contrast to Lie algebras such as SU($N$) and O($N$), which
are finite-dimensional, the Virasoro algebra is infinite-dimensional. (It does have a finite subalgebra, namely the one
generated by $L_1$, $L_0$, and $L_{-1}$.)

\subsubsection*{Lie groups and Lie algebras}

Let us recall some results from the theory of Lie groups and Lie algebras. For an element $U$ of the Lie group
SU($N$), one writes
\begin{equation}
  U = \exp(\ii\,\mathbf T\cdot\mathbf a) = \exp\left(\ii\sum_{j=1}^{o(G)}T_ja_j\right) ,
\end{equation}
where $T_j$ are the generators of the corresponding Lie algebra, satisfying
\begin{equation}\label{eq:cft-lie}
  [T_j,T_k] = \ii c_{jk\ell}T_\ell .
\end{equation}
Since $U^\dagger = U^{-1}$, the generators $T_j$ are hermitian, $T_j^\dagger = T_j$, and since $\det(U) = 1$, these
generators are traceless. The commutation relation \eqref{eq:cft-lie} is nonlinear, so the definition of the
structure constants $c_{jk\ell}$ requires a normalization of the generators. The standard normalization condition in
physics is $\Tr(T_jT_k) = (1/2)\delta_{jk}$.

\subsection{Highest-weight representations}

We next discuss representations of the Virasoro algebra. The term ``Verma module'' is used to describe these
representations. A representation is unitary if the states have non-negative norm (and a strictly positive norm for
nontrivial representations). One builds up the states in a representation via operation of Virasoro generators on a
highest-weight state.

\subsubsection*{The SU(2) algebra of angular momentum}

Recall how we do this in the case of the SU(2) algebra for angular momentum. The Lie algebra involves three
infinitesimal generators $J_j$, $j = 1,2,3$ obeying the commutation relation
\begin{equation}
  [J_j,J_k] = \ii\epsilon_{jk\ell}J_\ell ,
\end{equation}
so $c_{jk\ell} = \epsilon_{jk\ell}$. A specific form for the generators is $T_j = J_j = (1/2)\sigma_j$ (also denoted
$(1/2)\tau_j$), with $j = 1,2,3$, where the Pauli matrices $\sigma_j$ are
\begin{equation}
  \sigma_1 = \begin{pmatrix} 0 & 1\\ 1 & 0\end{pmatrix}, \qquad
  \sigma_2 = \begin{pmatrix} 0 & -\ii\\ \ii & 0\end{pmatrix}, \qquad
  \sigma_3 = \begin{pmatrix} 1 & 0\\ 0 & -1\end{pmatrix} .
\end{equation}
The operators $J_1$, $J_2$, and $J_3$ are often denoted as $J_x$, $J_y$, and $J_z$, especially in applications to the
physics of angular momenta. One defines raising and lowering operators in SU(2) as $J_\pm = J_1\pm\ii J_2$. The basic
SU(2) commutation relations are equivalent to
\begin{gather}
  [J_+,J_-] = 2J_3 \\
  [J_3,J_\pm] = \pm J_\pm .
\end{gather}
In an SU(2)-symmetric theory, the states form irreducible representations of the SU(2) algebra. These states are
labelled by $\ket{J,J_3} = \ket{j,m}$, where $m$ runs from $-j$ to $j$. The states are normalized according to
$\braket{jm}{j'm'} = \delta_{jj'}\delta_{mm'}$. With these states, the operators $J^2 = J_1^2+J_2^2+J_3^2$ and $J_3$
are diagonal and have eigenvalues $j(j+1)$ and $j_3 = j_z\equiv m$, respectively. The highest weight state is
$\ket{j,j}$. The operator $J_+$ is a raising operator, and its adjoint, $J_- = (J_+)^\dagger$, is a lowering operator.
The action of $J_+$ on the state $\ket{j,m}$ can be written as $J_+\ket{j,m} = a_{j,m,+}\ket{j,m'}$, where
$a_{j,m,\pm}$ is a constant depending on $j$, $m$, and the choice $\pm$. To see that $J_+$ is a raising operator, we
will show that $m' = m+1$. We have $[J_3,J_+] = J_+$, so $(J_3J_+-J_+J_3)\ket{j,m} = J_+\ket{j,m}$. Now, evaluating the
two terms on the left-hand side of this equation, we have $J_3J_+\ket{j,m} = a_{j,m,+}J_3\ket{j,m'}$ and
$J_+J_3\ket{j,m} = ma_{j,m,+}\ket{j,m'}$. The right-hand side is $a_{j,m,+}\ket{j,m'}$, so combining terms, we have
$(m'-m-1)a_{j,m,+}\ket{j,m'} = 0$ and hence $m' = m+1$. The corresponding steps show that $J_-$ is a lowering operator.
As shown in standard quantum mechanics textbooks, the $a_{j,m,\pm} = [(j\mp m)(j\pm m+1)]^{1/2}$. Thus,
\begin{align}
  J_+\ket{j,m} &= \sqrt{j(j+1)-m(m+1)}\,\ket{j,m+1} \nonumber\\
               &= \sqrt{(j-m)(j+m+1)}\,\ket{j,m+1} .
\end{align}
It follows that the highest-weight state $\ket{j,j}$ is annihilated by $J_+$:
\begin{equation}
  J_+\ket{j,j} = 0 ,
\end{equation}
and, correspondingly, $J_-\ket{j,-j} = 0$. The other states may be constructed by the operation of lowering operators
on this highest-weight state:
\begin{align}
  J_-\ket{j,m} &= \sqrt{j(j+1)-m(m-1)}\,\ket{j,m-1} \nonumber\\
               &= \sqrt{(j+m)(j-m+1)}\,\ket{j,m-1} .
\end{align}
Equivalently, one could start by calling the highest-weight state $\ket{j,-j}$, which is annihilated by $J_-$ and
constructing the other states via the operation of the raising operator $J_+$ on $\ket{j,-j}$. (As we recall, this $m$
can be integral or half-integral and should not be confused with an integer denoted $m$ used below in the discussion of
rational conformal algebras.)

\subsubsection*{The Verma module}

The construction of the states in a Verma module is similar but more complicated, for several reasons. First, the
algebra is infinite-dimensional, rather than finite-dimensional. Second, a unitary representation requires that a
state should have a non-negative norm, but this requires that one impose certain restrictions, as will be discussed
below. The analogue of $J_z\ket{j,m} = m\ket{j,m}$ is
\begin{equation}
  L_0\ket{c,h} = h\ket{c,h} ,
\end{equation}
where $h$ is the conformal dimension. From the basic commutation relation defining the Virasoro algebra, it follows
that $[L_0,L_n] = -nL_n$, so, applying this to a state $\ket{c,h}$, we have
$(L_0L_n-L_nL_0)\ket{c,h} = -nL_n\ket{c,h}$. Using $L_0\ket{c,h} = h\ket{c,h}$, we have
$(L_0L_n-hL_n+nL_n)\ket{c,h} = (L_0-h+n)L_n\ket{c,h} = 0$. With $L_n\ket{c,h}\propto\ket{c,h'}$, this implies
$(h'-h+n) = 0$, i.e., $h' = h-n$. This shows that $L_n$ is a lowering operator, which takes a state $\ket{c,h}$ to a
state $\ket{c,h-n}$. By similar calculations, it also follows that $L_{-n}$ is a raising operator, which takes a state
$\ket{c,h}$ to $\ket{c,h+n}$. Note that $L_{-n}$ can formally be regarded as the adjoint of $L_n$, i.e.,
$L_{-n} = (L_n)^\dagger$, just as $J_+ = (J_-)^\dagger$ in SU(2). One may define $\ket{h}$ to be a highest-weight
state, so that
\begin{equation}
  L_n\ket{h} = 0 \quad\text{if}\quad n\ge 1 .
\end{equation}
By convention, one may say that this state $\ket{h}$ is at level zero of the Verma module. Starting from the
highest-weight state, one can construct states at other levels by the operation of the raising operator. Thus, the
next state is $L_{-1}\ket{h}$, which we describe as being at level 1. We can get a level-2 state in two ways:
$L_{-2}\ket{h}$ and $(L_{-1})^2\ket{h}$, and we can get a level-3 state in three independent ways, $L_{-3}\ket{h}$,
$L_{-2}L_{-1}\ket{h}$, and $(L_{-1})^3\ket{h}$. (Note that $L_{-1}L_{-2}\ket{h}$ is not independent of
$L_{-2}L_{-1}\ket{h}$.) In general, the number of ways of constructing a state at level $\ell$ is equal to the number
of partitions of the (non-negative) integer $\ell$, denoted by the partition function, $p(\ell)$. (Note that the same
word, ``partition function'', is used here in the context of number theory, and should not be confused with its
different meaning in the context of statistical mechanics.) Simple examples of partitions of (positive) integers are
$p(0) = 1$, $p(1) = 1$,
\begin{align}
  p(2) &= 2:\ [(2),(1,1)]; \qquad p(3) = 3:\ [(3),(2,1),(1,1,1)]; \nonumber\\
  p(4) &= 5:\ [(4),(3,1),(2,2),(2,1,1),(1,1,1,1)] .
\end{align}

A state $\ket{s}$ at level $\ell = k_1+\cdots+k_m$ can be written as the result of operating on $\ket{h}$ with a set
of raising operators:
\begin{equation}
  \ket{s} = L_{-k_m}\cdots L_{-k_2}L_{-k_1}\ket{h} .
\end{equation}
This state is an eigenstate of $L_0$ with eigenvalue $h+\ell$: $L_0\ket{s} = (\ell+h)\ket{s}$. The complete set of
states $\ket{s}$ comprise the Verma module. One may define an inner product $\braket{s'}{s}$ on this space of states
using the property that $L_n = (L_{-n})^\dagger$. States at different levels $\ell$ and $\ell'\ne\ell$ are orthogonal.
The inner products of states at the same level are calculated by means of the Virasoro algebra, using the property
that $L_n\ket{h} = 0$ and thus also $\bra{h}L_{-n} = 0$. Just as the square of the norm of a vector $\mathbf v$ in a
vector space is given by $\lVert\mathbf v\rVert^2 = \braket{\mathbf v}{\mathbf v}$, so also the square of the norm of
this state, $\lVert\ket{s}\rVert^2$ in the Verma module is given by the inner product
\begin{equation}
  \lVert\ket{s}\rVert^2 = \braket{s}{s}
  = \bra{h}(L_{k_1}L_{k_2}\cdots L_{k_m})(L_{-k_m}\cdots L_{-k_2}L_{-k_1})\ket{h} .
\end{equation}

\subsection{Unitarity and the Kac determinant}

For example, the square of the norm of the state $L_{-n}\ket{h}$ is
\begin{align}
  \bra{h}L_nL_{-n}\ket{h} &= \bra{h}L_{-n}L_n + 2nL_0 + \frac{c}{12}n(n^2-1)\ket{h} \nonumber\\
  &= \left[2nh + \frac{c}{12}n(n^2-1)\right]\braket{h}{h} .
\end{align}
If $c<0$, then this is negative for sufficiently large $n$, so $c$ must be non-negative for a unitary representation,
i.e., $c\ge0$. Furthermore, setting $n = 1$, one sees that the condition that the norm be non-negative requires that
$h\ge0$.

\subsubsection*{The Gram matrix}

The determination of restrictions due to the requirement of unitarity consists of studying the eigenvalues of the
matrices whose entries are the inner products of all states at a given level $\ell$ (since states at different levels
are orthogonal),
\begin{equation}
  M_{ij}^{(\ell)} = \braket{s_i}{s_j} .
\end{equation}
This is called the Gram matrix. One often writes this matrix in terms of the basic norm $\braket{h}{h}$. Since (for
real $c$ and $h$) $M^{(\ell)}$ is a real symmetric matrix, it can be diagonalized by an orthogonal operator, yielding a
set of $p(\ell)$ eigenstates and associated eigenvalues, which are the squares of the norms of the eigenstates. Since
$(M^{(\ell)})^\dagger = (M^{(\ell)})^T = M^{(\ell)}$, it follows that the eigenvalues are real,
but this does not imply that they are non-negative. The requirement that the eigenvalues be non-negative yields further
restrictions, as discussed next. The trace and determinant of $M^{(\ell)}$ are equal to the sum and product of these
eigenvalues:
\begin{equation}
  \Tr(M^{(\ell)}) = \sum_{j=1}^{p(\ell)}\lambda_j^{(\ell)} , \qquad
  \det(M^{(\ell)}) = \prod_{j=1}^{p(\ell)}\lambda_j^{(\ell)} .
\end{equation}
For each level $\ell$, a necessary condition for the representation to be unitary is that the trace and determinant
must be non-negative. At level $\ell = 0$, we have $M^{(0)} = \braket{h}{h}$. It is convenient to divide by this
overall factor when discussing matrices of states for higher levels, and we shall follow this convention. Then at
level 1,
\begin{equation}
  M^{(1)} = \bra{h}L_1L_{-1}\ket{h} = 2\bra{h}L_0\ket{h} = 2h .
\end{equation}

\subsubsection*{Level 2}

At level 2, since $p(2) = 2$, there are two states in the Verma module, and correspondingly a $2\times2$ matrix of
inner products. We define the entries as follows:
\begin{align}
  M_{11}^{(2)} &= \bra{h}L_2L_{-2}\ket{h} \\
  M_{12}^{(2)} &= \bra{h}L_2(L_{-1})^2\ket{h} \\
  M_{21}^{(2)} &= \bra{h}(L_1)^2L_{-2}\ket{h} = M_{12}^{(2)} \\
  M_{22}^{(2)} &= \bra{h}(L_1)^2(L_{-1})^2\ket{h} .
\end{align}
Evaluating these inner products, one has
\begin{equation}
  M^{(2)} = \begin{pmatrix} 4h+(c/2) & 6h\\ 6h & 4h(2h+1)\end{pmatrix} .
\end{equation}
As a special case of the general restrictions, the trace and determinant must be non-negative for a unitary
representation. We have
\begin{equation}
  \Tr(M^{(2)}) = 8h(h+1) + \frac{c}{2}
\end{equation}
and
\begin{align}
  \det(M^{(2)}) &= 2h[16h^2+2(c-5)h+c] \nonumber\\
  &= 32(h-h_{1,1})(h-h_{2,1})(h-h_{1,2}) ,
\end{align}
where
\begin{equation}
  h_{1,1} = 0
\end{equation}
and
\begin{equation}
  \{h_{2,1},h_{1,2}\} = \frac{1}{16}\left[5-c\pm\sqrt{(1-c)(25-c)}\right] .
\end{equation}
We consider $c$ in the interval $0\le c\le1$. Note that if $c = 1$, then $h_{2,1} = h_{1,2} = 1/4$; as $c$ decreases
from 1 toward 0, $h_{2,1}$ and $h_{1,2}$ move along an upper and lower branch of a parabola which is concave toward the
left in a diagram with horizontal axis $c$ and vertical axis $h$. These two branches approach the limits
$h_{2,1} = 5/8$ and $h_{1,2} = 0$ as $c\to0$ from above. The determinant is negative in the region $0<c<1$ bounded
above and below by the upper and lower branches of this parabolic curve. This is the region interior to the parabolic
curve. This means that one of the two eigenvalues of $M^{(2)}$ is negative and hence these values of $(c,h)$ do not
yield a unitary representation of the Virasoro algebra. This can be seen from an examination of the eigenvalues, which
are
\begin{equation}
  \lambda_{1,2}^{(2)} = 4h(h+1) + \frac{c}{4} \pm \frac{1}{4}\left[256h^4+576h^2-32ch^2+c^2\right]^{1/2} .
\end{equation}

\subsubsection*{The Kac determinant and the unitary series}

One continues in this way to analyze the properties of the sets of states at higher levels $\ell\ge3$. A general
mathematical result due to V.~Kac is a formula for the determinant of $M^{(\ell)}$, namely
\begin{equation}\label{eq:cft-kac}
  \det(M^{(\ell)}) = \alpha_\ell\prod_{r,s\ge1}[h-h_{r,s}]^{p(\ell-rs)}
\end{equation}
where $1\le r,\,s\le\ell$ subject to the condition that $rs\le\ell$. In Eq.~\eqref{eq:cft-kac}, $\alpha_\ell$ is a
positive numerical prefactor that is independent of $c$ and $h$ and hence does not affect the analysis of how the
signs of the eigenvalues depend on $c$ and $h$.

It turns out that in the interval $0\le c<1$, the solutions for $c$ and the corresponding scaling dimensions $h_{r,s}$
that yield unitary conformal algebras are the following discrete set:
\begin{equation}
  c = 1-\frac{6}{m(m+1)}, \qquad m\ge2 \text{ (integral)} ,
\end{equation}
with the scaling dimensions of operators
\begin{equation}\label{eq:cft-hrs}
  h_{r,s} = \frac{[(m+1)r-ms]^2-1}{4m(m+1)}, \qquad\text{with}\quad 1\le s\le r\le m-1 ,
\end{equation}
and similarly with $\bar h_{r,s}$ in the anti-holomorphic Virasoro algebra of the $\bar L_n$. The $h_{r,s}$ can also be
expressed directly in terms of $c$ according to the relations
\begin{equation}
  h_{r,s} = \frac{c-1}{24} + \frac{(ra_++sa_-)^2}{4} ,
\end{equation}
where
\begin{equation}
  a_\pm = \frac{\sqrt{1-c}\pm\sqrt{25-c}}{\sqrt{24}} .
\end{equation}
(See Di Francesco, Mathieu, and S\'en\'echal, \emph{Conformal Field Theory}, pp.~207--209, or the other references
listed at the end.) The case $m = 2$ and hence $c = 0$ is trivial, since in this case the only values of $r$ and $s$
are $r = s = 1$, so that there is only one scaling dimension, namely $h_{1,1} = 0$, corresponding to the identity
operator. More generally, for any $c$, it follows that $h_{1,1} = 0$, again corresponding to the identity operator,
which is automatically an element of any conformal algebra. Note that $h_{r,s}$ is invariant under the simultaneous
transformations $r\to m-r$ and $s\to m+1-s$, i.e.,
\begin{equation}\label{eq:cft-hsym}
  h_{r,s} = h_{m-r,m+1-s} ,
\end{equation}
so that one may restrict $r$ and $s$ to the ranges
\begin{equation}
  h_{r,s}:\ 1\le r\le m-1, \qquad 1\le s\le r .
\end{equation}
For this unitary series, the integer $m$ is related to $c$ according to
\begin{equation}\label{eq:cft-m}
  m = \frac{1}{2}\left[-1\pm\left(\frac{25-c}{1-c}\right)^{1/2}\right] .
\end{equation}
The two signs in front of the square root in Eq.~\eqref{eq:cft-m} are equivalent, since reversing the sign
interchanges $m$ with $-(m+1)$, which amounts to an interchange of $r$ and $s$ in Eq.~\eqref{eq:cft-hrs}. This set of
theories comprises what are called rational (unitary) 2D conformal field theories, corresponding to the given
conformal algebras. The correspondence of the unitary conformal algebras with statistical mechanical models was
discussed by A.~A.~Belavin, A.~M.~Polyakov and A.~B.~Zamolodchikov (BPZ), Nucl.\ Phys.\ B \textbf{241}, 333 (1984);
J.\ Stat.\ Phys.\ \textbf{34}, 763 (1984); for reviews, see the references at the end.

\subsection{The 2D Ising model as a conformal field theory}

An important special case is $m = 3$, corresponding to $c = 1/2$, which is the 2D Ising ferromagnet. Here,
$1\le s\le r\le2$, so we have
\begin{equation}
  h_{r,s} = \{h_{1,1},\ h_{2,1},\ h_{2,2}\} .
\end{equation}
From Eq.~\eqref{eq:cft-hrs},
\begin{equation}
  h_{1,1} = 0, \qquad h_{2,1} = \frac{1}{2}, \qquad h_{2,2} = \frac{1}{16} ,
\end{equation}
and similarly for the $\bar h_{r,s}$. From the symmetry relation \eqref{eq:cft-hsym}, it follows that
$h_{1,2} = h_{2,2}$. The actual scaling dimensions here are thus
\begin{equation}
  x_{r,s} = h_{r,s}+\bar h_{r,s} = 2h_{r,s} .
\end{equation}
This gives the values
\begin{equation}
  x_{1,1} = 0, \qquad x_{2,1} = 1, \qquad x_{2,2} = \frac{1}{8} .
\end{equation}
To construct a conformal field theory corresponding to this conformal algebra, one then constructs a set of primary
operators which, operating on a vacuum state, produce the CFT states such that $L_0\ket{h} = h\ket{h}$. As noted, this
always contains the identity operator, 1. In general, if an operator $\Order$ has dimension $x_\Order$, then its
connected two-point function is
\begin{equation}
  \ev{\Order(0)\Order(r)} = \frac{1}{|r|^{2x_\Order}} .
\end{equation}
Now
\begin{equation}
  \ev{\sigma(0)\sigma(r)} \sim \frac{1}{r^{2x_\sigma}} = \frac{1}{r^{d-2+\eta}} = \frac{1}{r^\eta} ,
\end{equation}
so since $\eta = 1/4$, we identify the (primary) operator $\sigma$ with the scaling dimension
\begin{equation}
  x_\sigma = x_{2,2} = 2h_{2,2} = \frac{1}{8} .
\end{equation}
Furthermore, a property of the 2D Ising model at criticality is that, with $\sigma_i\sigma_j\sim\epsilon$, an energy
operator (do not confuse with the use of the symbol $\epsilon = 4-d$ in the $\epsilon$ expansion), one has
\begin{equation}
  \ev{\epsilon(0)\epsilon(r)} \sim \frac{1}{r^2} .
\end{equation}
So this primary operator $\epsilon$ has scaling dimension
\begin{equation}
  x_\epsilon = x_{2,1} = 2h_{2,1} = 1 .
\end{equation}

\subsubsection*{Fusion rules and the perturbed theory}

The CFT involves an operator algebra (also called the fusion rules, example of very simple operator product
expansions)
\begin{gather}
  \sigma\times\sigma = 1+\epsilon \\
  \sigma\times\epsilon = \sigma \\
  \epsilon\times\epsilon = 1
\end{gather}

An easy way to see the dimensions of primary operators is to generalize the Ising model to a Euclidean quantum field
theory invariant under the same $Z_2$ global symmetry. This is the perturbed conformal field theory with generating
function
\begin{equation}
  Z = \int\prod_s\dd\sigma\,e^{-S} ,
\end{equation}
where the action $S$ is (to use the notation standard in CFT)
\begin{equation}
  S_{\mathrm{IFT}} = S_{c=1/2} + \tau\int\dd^2x\,\epsilon(x) + h\int\dd^2x\,\sigma(x) ,
\end{equation}
where $S_{c=1/2}$ stands for the action of the conformal field theory (CFT) with central charge $c = 1/2$, and
$\tau = (K-K_c)/K_c$. The terms proportional to $\tau$ and $h$ in the action are, respectively, temperature-like and
magnetic-field-like perturbations of the critical conformal $c = 1/2$ theory. These two terms are the two relevant
perturbations, involving the deviation of the inverse temperature from $K_c$, coupled to $\epsilon$, and the external
magnetic field $h$ (i.e., the deviation from its critical value of zero) coupled to the field, $\sigma$. Here, the
primary fields $\sigma$ and $\epsilon$ have conformal dimensions $x_\sigma = 1/8$ and $x_\epsilon = 1$. The conformal
normalizations of these fields are $D_\epsilon = 1$, $D_\sigma = 1/8$, in accordance with
\begin{equation}
  \ev{\epsilon(0)\epsilon(r)} \sim \frac{1}{r^2}
\end{equation}
and
\begin{equation}
  \ev{\sigma(0)\sigma(r)} \sim \frac{1}{r^{1/4}} .
\end{equation}
We can then relate these conformal dimensions to the critical exponents that we have discussed earlier. Since the
action is dimensionless (in units where $\hbar = 1$), for the first term, we have
$0 = y_\tau-2+D_\epsilon = 1-2+1$, and for the second term we have $0 = y_h-2+D_\sigma = 15/8-2+1/8$.

This, then, summarizes the analysis for the Ising case, $m = 3$, for which the central charge is $c = 1/2$. In a
similar manner, one can associate conformal algebras with other statistical mechanical models:
\begin{enumerate}[nosep]
  \item $m = 4$, $c = 7/10 = 0.70$: tricritical Ising
  \item $m = 5$, $c = 4/5 = 0.80$: $q = 3$ Potts
\end{enumerate}
etc. Unitary representations also exist for $c\ge1$ (and $h\ge0$).

\subsubsection*{Further reading}

Some references for further reading on conformal field theory include
\begin{enumerate}[nosep]
  \item A.~A.~Belavin, A.~M.~Polyakov, and A.~B.~Zamolodchikov, op.\ cit.\ (1984).
  \item P.~Di Francesco, P.~Mathieu, and D.~S\'en\'echal, \emph{Conformal Field Theory} (Springer, New York, 1996).
  \item C.~Itzykson, H.~Saleur, and J.-B.~Zuber, \emph{Conformal Invariance and Applications to Statistical
    Mechanics} (World Scientific, Singapore, 1988).
  \item V.~G.~Kac and A.~K.~Raina, \emph{Highest-Weight Representations of Infinite Dimensional Lie Algebras}
    (World Scientific, Singapore, 1987).
  \item P.~Ginsparg, \emph{Applied Conformal Field Theory} (Les Houches lecture notes).
  \item A.~N.~Schellekens, \emph{Conformal Field Theory} (Lecture notes).
\end{enumerate}

Many of the advances associated with the applications of conformal algebra to critical properties of second-order
phase transitions in 2D occurred in 1983 and throughout the 1980s. For an earlier review of the renormalization group,
see M.~E.~Fisher, Rev.\ Mod.\ Phys.\ \textbf{46}, 597 (1974).

More recently, there has been another advance using conformal field theory in higher dimensions. Here, the conformal
group is finite-dimensional rather than infinite-dimensional, so the methods used are different. For a review, see,
e.g., D.~Poland and D.~Simmons-Duffin, Nature Physics \textbf{12}, 535 (2016) and papers on the arXiv by Poland,
Simmons-Duffin, Rastelli, and others. A recent paper relating this to gauge theories is Z.~Li and D.~Poland, JHEP 03
(2021) 172 [arXiv:2005.01721].
